% Regression: rv4061_ex_grid-benchmarks — renders Gauge equation B^i = X A^i X^{-1} between D x D matrices: both…
% Formula: Gauge equation B^i = X A^i X^{-1} between D x D matrices: both sides
\documentclass[varwidth=16cm, border=6pt]{standalone}
\usepackage{amsmath, amssymb}
\usepackage{tenkz}
\begin{document}

B1:
\begin{tenkz}[periodic, physical=up, bond label={$D$ at 1-2}]
  \tn[up=$i_1$]{A} & \tn[up=$i_2$]{A} & \tn[up=$i_3$]{A}
\end{tenkz}

\medskip
B2:
% Gauge equation B^i = X A^i X^{-1} between D x D matrices: both sides
% carry the same boundary signature (2 open virtual + 1 physical); the
% protruding stubs ARE the matrix indices.
\[
  \begin{tenkz}[physical=up]
    \tn[up=$i$]{B}
  \end{tenkz}
  \;=\;
  \begin{tenkz}[physical=up]
    \tnX{X} & \tn[up=$i$]{A} & \tnX{X^{-1}}
  \end{tenkz}
\]

B3:
% E_A = sum_i A^i (x) conj(A^i): a map on the doubled virtual space --
% four open virtual legs; west pair = input (X plugs in there), east pair
% = output.  The contracted physical leg IS the sum over i (contrast B8,
% where the sum is explicit and the leg stays open).
\[
  \mathcal{E}_A \;=\;
  \begin{tenkz}[sandwich]
    \tn{A} \\
    \tn*{A}
  \end{tenkz}
\]

B4:
% TI blocking: A^{(L)}_{(i_1...i_L)} = A^{i_1} ... A^{i_L} as D x D
% matrices -- the virtual boundary indices stay open on BOTH sides; the
% physical counts differ only by grouping (L legs vs one combined index).
\[
  \begin{tenkz}[physical=up]
    \tn[up=$i_1$]{A}\tnspan[brace below]{4}{L\text{ copies}} &
    \tn[up=$i_2$]{A} & \tndots & \tn[up=$i_L$]{A}
  \end{tenkz}
  \;=\;
  \begin{tenkz}[physical=up]
    \tn[pill, wide=2, up={$i_1\cdots i_L$}]{A^{(L)}}
  \end{tenkz}
\]

B7:
% Zipper (pulling through, arXiv:2203.12563): V (O-window (x) A-window)
%   = (B-window) V, with V : C^{D_O} x C^{D_A} -> C^{D_B} fusing the two
% bond spaces and B^i = sum_j V (O^{i j} (x) A^j) the fused site.  LHS:
% V's single combined leg is the open west index; the east O/A pair stays
% open.  RHS: the fused-bond B word; the bond-SPLITTING glyph (V at the
% east end of a one-row word) is a tracked Phase-1 primitive, so V closes
% the equation as a symbol here.
\[
  \begin{tenkz}[rows={op, ket}]
    \tnfuse[span=2]{V} & \tn[mpo]{O}\tnspan[brace above]{3}{n} & \tndots & \tn[mpo]{O} \\
                       & \tn{A}                                & \tndots & \tn{A}
  \end{tenkz}
  \;=\;
  \begin{tenkz}[physical=up]
    \tn[box, up=$i_1$]{B}\tnspan[brace above]{3}{n} & \tndots & \tn[box, up=$i_n$]{B}
  \end{tenkz}
  \,V
\]

% E_A(X) = sum_i A^i X (A^i)^dagger: the sum index i is explicit, so no
% contracted physical leg is drawn -- each summand is a product of three
% matrices on ONE virtual wire, open at both ends (a map, not a scalar).
B8 inline: since
$\mathcal{E}_A(X)=\sum_{i}\,\tnpic[inline]{\tn[box]{A^i} & \tnX{X} & \tn*[box]{A^i}}$
each summand is a congruence.

\end{document}
